\chapter{Przestrzenie polskie i zbiory borelowskie}
\label{ch:topog-polskie-borelowskie}

\noindent\textbf{Cel rozdziału.}
Zupełność metryczna z rozdziału~11 i metryzowalność
z rozdziału~10 prowadzą do pojęcia przestrzeni
polskiej. Zbadamy jej trwałość przy podstawowych
konstrukcjach oraz scharakteryzujemy te podprzestrzenie
przestrzeni polskich, które same są polskie.
Następnie zbudujemy $\sigma$-algebrę zbiorów
borelowskich oraz jej hierarchię.
Od rozdziału~20 oddzielimy bardziej zaawansowane
zagadnienia opisowej teorii mnogości, w szczególności
zbiory analityczne i twierdzenia o ich strukturze.

\section{Przestrzenie polskie: definicja i przykłady}
\label{sec:topog-polskie-definicja}

\begin{definicja}[Przestrzeń polska]
\label{def:topog-polskie}
Przestrzeń topologiczną $X$ nazywamy
\emph{przestrzenią polską}, jeśli istnieje
metryka $d$ wyznaczająca jej topologię,
taka, że $(X,d)$ jest zupełna,
oraz $X$ ma przeliczalny zbiór gęsty.
Nie wymagamy, aby każda metryka
zgodna z topologią $X$ była zupełna.
\end{definicja}

\begin{stwierdzenie}[Drugi aksjomat przeliczalności]
\label{prop:topog-polskie-drugi-aksjomat}
Każda przestrzeń polska spełnia
drugi aksjomat przeliczalności
i jest przestrzenią Hausdorffa.
\end{stwierdzenie}
\begin{proof}
Wybierzmy zgodną metrykę $d$
i przeliczalny zbiór gęsty $D\subseteq X$.
Rodzina
$\{B_d(a,r):a\in D,\ r\in\mathbb Q,\ r>0\}$
jest przeliczalna. Jest bazą:
jeśli $x\in U$ i $U$ jest otwarty,
istnieje $\varepsilon>0$ z
$B_d(x,\varepsilon)\subseteq U$.
Wybierzmy $a\in D$ z
$d(x,a)<\varepsilon/2$
oraz wymierne $r$ z
$d(x,a)<r<\varepsilon-d(x,a)$.
Wówczas
$x\in B_d(a,r)\subseteq U$.
Hausdorffowość wynika z metryzowalności.
\end{proof}

\begin{przyklad}[Znane przestrzenie polskie]
\label{ex:topog-polskie-podstawowe}
Przestrzeń $\mathbb R^m$ ze zwykłą
metryką euklidesową jest zupełna,
a $\mathbb Q^m$ jest gęstym zbiorem
przeliczalnym. Każda jej domknięta
podprzestrzeń, z metryką ograniczoną,
jest również zupełna i ośrodkowa.

Przedział $(0,1)$ jest polski,
chociaż metryka euklidesowa
ograniczona do $(0,1)$ nie jest
zupełna. Zgodną metrykę zupełną
można otrzymać przez zanurzenie
$t\mapsto\log(t/(1-t))$
w prostą: odwzorowanie to jest
homeomorfizmem $(0,1)$ na $\mathbb R$,
a więc metryka
\[
d_\ast(s,t)=
\left|\log\frac{s}{1-s}
-\log\frac{t}{1-t}\right|
\]
jest zgodna i zupełna.
\end{przyklad}

\begin{przyklad}[Przestrzenie Cantora i Baire'a]
\label{ex:topog-polskie-cantor-baire}
Niech $A$ oznacza odpowiednio
$\{0,1\}$ lub $\mathbb N$,
w obu przypadkach z topologią dyskretną.
Na zbiorze $A^{\mathbb N}$ ciągów
określmy metrykę
\[
d_A(x,y)=
\begin{cases}
0,&x=y,\\
2^{-k},&k=\min\{n\ge1:x_n\ne y_n\}.
\end{cases}
\]
Jest to ultrametryka: jeżeli
$x_k=y_k$ i $y_k=z_k$
dla każdego $k<m$, to również
$x_k=z_k$ dla tych indeksów;
stąd
$d_A(x,z)\le\max\{d_A(x,y),d_A(y,z)\}$.
Kule o odpowiednio małym promieniu
wyznaczają dokładnie warunek
zgodności skończonej liczby
początkowych współrzędnych,
zatem $d_A$ wyznacza topologię
produktową.

Przestrzeń ta jest zupełna.
Jeżeli $(x^{(j)})_{j\ge1}$ jest
ciągiem Cauchy'ego, to dla każdego
$n$ współrzędne
$x_1^{(j)},\ldots,x_n^{(j)}$
są ostatecznie stałe:
wystarczy wykorzystać oszacowanie
$d_A(x^{(j)},x^{(l)})<2^{-n}$.
Dla każdej współrzędnej $k$
jej ostateczna wartość jest
jednoznacznie określona.
Wartości te definiują
$x\in A^{\mathbb N}$.
Dla ustalonego $n$ wyrazy
ciągu od pewnego miejsca
zgadzają się z $x$
na pierwszych $n$ współrzędnych,
zatem ich odległość od $x$
jest nie większa niż
$2^{-(n+1)}$.
Ponieważ $n$ jest dowolne,
otrzymujemy zbieżność do $x$.

Ciągi, których wszystkie współrzędne
od pewnego miejsca są równe
ustalonemu $a_0\in A$,
tworzą zbiór przeliczalny i gęsty.
Dla $A=\{0,1\}$ otrzymujemy
\emph{przestrzeń Cantora}
$2^{\mathbb N}$, a dla
$A=\mathbb N$ ---
\emph{przestrzeń Baire'a}
$\mathbb N^{\mathbb N}$.
Obie są polskie; pierwsza jest
zwarta na mocy twierdzenia Tychonoffa,
a druga nie jest zwarta, ponieważ
jej obraz przez rzutowanie na
pierwszą współrzędną jest
nieskończoną dyskretną przestrzenią
$\mathbb N$, która nie jest zwarta.
\end{przyklad}

\begin{twierdzenie}[Przeliczalne iloczyny przestrzeni polskich]
\label{thm:topog-polskie-iloczyn}
Niech $(X_n)_{n\ge1}$ będzie ciągiem
przestrzeni polskich.
Wtedy przestrzeń
$X=\prod_{n\ge1}X_n$
z topologią produktową jest polska.
\end{twierdzenie}
\begin{proof}
Jeśli któryś czynnik jest pusty,
to $X=\varnothing$, a przestrzeń
pusta jest polska. Załóżmy więc,
że wszystkie $X_n$ są niepuste.
Wybierzmy na $X_n$ zgodną
zupełną metrykę $d_n$ i
przeliczalny zbiór gęsty $D_n$.
Dla $x=(x_n)$ oraz $y=(y_n)$
określmy
\[
d(x,y)=\sum_{n=1}^{\infty}
2^{-n}\min\{1,d_n(x_n,y_n)\}.
\]
Nierówność trójkąta wynika z
$\min\{1,a+b\}\le
\min\{1,a\}+\min\{1,b\}$,
a dodatniość metryki z faktu,
że różne punkty mają różną
co najmniej jedną współrzędną.
Szereg jest zbieżny, gdyż każdy
składnik jest ograniczony przez $2^{-n}$.

Topologia metryki $d$ jest produktowa.
Zbieżność w $d$ wymusza zbieżność
każdej współrzędnej, ponieważ
$2^{-n}\min\{1,d_n(x_n,y_n)\}\le d(x,y)$.
W drugą stronę skończenie wiele
współrzędnych można kontrolować
równocześnie, a ogon szeregu
$\sum_{n>N}2^{-n}$ uczynić
dowolnie małym. To dowodzi
zgodności topologii.

Jeśli $(x^{(j)})$ jest ciągiem
Cauchy'ego w $d$, każda
$(x_n^{(j)})$ jest ciągiem
Cauchy'ego w $d_n$, więc
ma granicę $x_n\in X_n$.
Oszacowanie ogona szeregu
pokazuje następnie, że
$x^{(j)}\to(x_n)$ w $d$.
Zatem $d$ jest zupełna.

Wybierzmy punkt bazowy
$a=(a_n)\in X$.
Zbiór ciągów różniących się od
$a$ tylko na skończenie wielu
współrzędnych, przyjmujących
tam wartości z odpowiednich $D_n$,
jest przeliczalny jako
przeliczalna suma skończonych
iloczynów zbiorów przeliczalnych.
Jest gęsty, ponieważ bazowy
zbiór otwarty topologii produktowej
ogranicza tylko skończenie wiele
współrzędnych. To kończy dowód.
\end{proof}

\section{Podprzestrzenie polskie i zbiory typu \texorpdfstring{$G_\delta$}{G-delta}}
\label{sec:topog-polskie-gdelta}

\begin{definicja}[Zbiory typu F-sigma i G-delta]
\label{def:topog-polskie-fsigma-gdelta}
Niech $X$ będzie przestrzenią
topologiczną. Zbiór $A\subseteq X$
jest typu $F_\sigma$, jeżeli
istnieją zbiory domknięte
$F_n\subseteq X$ takie, że
$A=\bigcup_{n\ge1}F_n$.
Jest typu $G_\delta$, jeżeli
istnieją otwarte $U_n\subseteq X$
takie, że $A=\bigcap_{n\ge1}U_n$.
\end{definicja}

\begin{lemat}[Zbiory otwarte i domknięte w przestrzeni metrycznej]
\label{lem:topog-polskie-otwarte-fsigma}
Niech $(X,d)$ będzie przestrzenią
metryczną. Każdy jej zbiór otwarty
jest typu $F_\sigma$, a każdy
zbiór domknięty jest typu $G_\delta$.
\end{lemat}
\begin{proof}
Niech $U\subseteq X$ będzie otwarty
oraz $F=X\setminus U$.
Jeśli $F=\varnothing$, mamy
$U=X$, który jest domknięty.
W przeciwnym przypadku
dla $n\ge1$ połóżmy
$F_n=\{x\in X:d(x,F)\ge1/n\}$.
Funkcja $x\mapsto d(x,F)$
jest ciągła, zatem $F_n$ jest
domknięty. Ponieważ dla
$x\in U$ istnieje kula
$B_d(x,\varepsilon)\subseteq U$,
mamy $d(x,F)>0$ i punkt $x$
należy do pewnego $F_n$.
Jeżeli $x\in F$, to
$d(x,F)=0$ i $x$ nie należy
do żadnego $F_n$. Stąd
$U=\bigcup_{n\ge1}F_n$.
Dla domkniętego $C$ stosujemy
pierwszą część do otwartego
$X\setminus C$ i bierzemy
dopełnienia.
\end{proof}

\begin{twierdzenie}[Charakterystyka polskich podprzestrzeni]
\label{thm:topog-polskie-gdelta}
Niech $X$ będzie przestrzenią polską
i $A\subseteq X$. Podprzestrzeń
$A$ jest polska wtedy i tylko wtedy,
gdy $A$ jest zbiorem typu $G_\delta$
w $X$.
\end{twierdzenie}
\begin{proof}
Niech najpierw
$A=\bigcap_{n\ge1}U_n$,
gdzie $U_n$ są otwarte w $X$.
Wybierzmy zupełną zgodną metrykę
$d$ na $X$. Gdy $A=\varnothing$,
twierdzenie jest oczywiste.
Dla każdego $n$ z
$X\setminus U_n\ne\varnothing$
określamy ciągłą na $A$ funkcję
$h_n(x)=1/d(x,X\setminus U_n)$.
Gdy $U_n=X$, przyjmujemy $h_n=0$.
Mianowniki są dodatnie dla
$x\in A$, ponieważ $U_n$ jest
otwarty i $x\in U_n$.
Na $A$ określmy
\[
\rho(x,y)=d(x,y)+
\sum_{n\ge1}2^{-n}
\min\{1,|h_n(x)-h_n(y)|\}.
\]
Jest to metryka: składniki
są pseudometrykami, a pierwszy
jest metryką. Wyznacza
topologię podprzestrzeni $A$.
Istotnie, $\rho\ge d$, zatem
zbieżność w $\rho$ implikuje
zbieżność w $d$; odwrotnie,
ciągłość każdego $h_n$ i
jednostajnie mały ogon szeregu
zapewniają zbieżność w $\rho$
przy zbieżności w $d$:
dla ustalonego $\varepsilon>0$
wybieramy $N$ z
$\sum_{n>N}2^{-n}<\varepsilon/2$,
a następnie kontrolujemy
skończenie wiele różnic
$|h_n(x_j)-h_n(x)|$, $n\le N$.
W przestrzeni metrycznej
zgodność zbiegających ciągów
do wszystkich punktów wystarcza
do równości topologii.

Udowodnimy zupełność.
Niech $(x_j)$ będzie ciągiem
Cauchy'ego w $\rho$.
Wówczas jest ciągiem Cauchy'ego
w $d$, a więc $x_j\to x$
dla pewnego $x\in X$.
Dla każdego ustalonego $n$
ciąg liczb $(h_n(x_j))$
jest ciągiem Cauchy'ego.
Istotnie, dla
$0<\varepsilon<1$ wybieramy
$j_0$ tak, by
$\rho(x_j,x_k)<2^{-n}\varepsilon$
dla wszystkich $j,k\ge j_0$.
Wtedy
$\min\{1,|h_n(x_j)-h_n(x_k)|\}
<\varepsilon<1$,
a więc
$|h_n(x_j)-h_n(x_k)|<\varepsilon$.
To wystarcza również dla
większych wartości $\varepsilon$.
Jeśli $x\notin U_n$,
to $d(x_j,X\setminus U_n)
\le d(x_j,x)\to0$, przez co
$h_n(x_j)\to+\infty$,
wbrew własności Cauchy'ego
(gdy $U_n=X$, taka sytuacja
nie występuje).
Zatem $x\in\bigcap_n U_n=A$.
Dla każdego $n$ ciągłość
$h_n$ na $A$ daje
$h_n(x_j)\to h_n(x)$,
ponieważ $x_j\to x$ w $d$
i wykazaliśmy $x\in A$.
Dla zadanego $\varepsilon>0$
dobieramy $N$ z
$\sum_{n>N}2^{-n}<\varepsilon/3$,
a następnie $j_0$ tak, by
$d(x_j,x)<\varepsilon/3$
oraz
$\sum_{n=1}^N2^{-n}
\min\{1,|h_n(x_j)-h_n(x)|\}
<\varepsilon/3$
dla $j\ge j_0$.
Wówczas $\rho(x_j,x)<\varepsilon$.
Zatem $(A,\rho)$ jest zupełna.
Podprzestrzeń przestrzeni ośrodkowej
metrycznej jest ośrodkowa,
ponieważ spełnia drugi aksjomat
przeliczalności jako podprzestrzeń
przestrzeni z przeliczalną bazą,
a wybór po jednym punkcie z
każdego niepustego zbioru tej bazy
daje przeliczalny zbiór gęsty.
Stąd $A$ jest polska.

Odwrotnie, załóżmy, że
$A$ jest polska.
Wybierzmy zgodną metrykę $d$
na $X$ i zgodną zupełną
metrykę $\sigma$ na $A$.
Dla każdego $a\in A$ i $n\ge1$
istnieje $r_n(a)>0$ takie, że
\[
r_n(a)<2^{-n},\qquad
\operatorname{diam}_{\sigma}
\bigl(B_d(a,2r_n(a))\cap A\bigr)
<2^{-n}.
\]
Istnienie wynika z ciągłości
identyczności
$(A,d|_A)\to(A,\sigma)$:
można wybrać tak małą kulę,
aby wszystkie jej punkty z $A$
leżały w kuli
$B_\sigma(a,2^{-n-2})$.
Niech
\[
V_n=\bigcup_{a\in A}
 B_d(a,r_n(a)).
\]
Są to zbiory otwarte w $X$
i $A\subseteq\bigcap_n V_n$.

Pokażemy inkluzję odwrotną.
Ustalmy $x\in\bigcap_nV_n$.
Dla każdego $n$ wybierzmy
$a_n\in A$ z
$d(x,a_n)<r_n(a_n)$.
Ponieważ $r_n(a_n)<2^{-n}$,
mamy $a_n\to x$ w metryce $d$.
Ustalmy $n$ oraz połóżmy
$\delta_n=r_n(a_n)-d(x,a_n)>0$.
Dla wystarczająco dużych $m$
mamy $r_m(a_m)<\delta_n$,
a więc
\[
d(a_m,a_n)
\le d(a_m,x)+d(x,a_n)
<r_m(a_m)+d(x,a_n)
<r_n(a_n).
\]
Wszystkie dostatecznie późne
$a_m$ leżą zatem w
$B_d(a_n,2r_n(a_n))\cap A$,
którego średnica w metryce
$\sigma$ jest mniejsza niż
$2^{-n}$. Ustalmy teraz dowolne
$\varepsilon>0$ i wybierzmy
$n$ takie, że $2^{-n}<\varepsilon$.
Z poprzedniego oszacowania
istnieje $m_0$ takie, że
dla wszystkich $m\ge m_0$
mamy
$a_m\in B_d(a_n,2r_n(a_n))\cap A$.
Zatem dla dowolnych
$m,l\ge m_0$ zachodzi
$\sigma(a_m,a_l)<2^{-n}<\varepsilon$.
To dowodzi, że $(a_m)$
jest ciągiem Cauchy'ego w $\sigma$. Z zupełności
otrzymujemy granicę $a\in A$.
Metryki $\sigma$ i $d|_A$
wyznaczają tę samą topologię,
więc $a_m\to a$ również w $d$.
Jednoznaczność granicy w
metrycznej $X$ i zbieżność
$a_m\to x$ dają $x=a\in A$.
Zatem $A=\bigcap_nV_n$.
\end{proof}

\begin{przyklad}[Liczby wymierne i niewymierne]
\label{ex:topog-polskie-wymierne}
Zbiór $\mathbb Q\subseteq\mathbb R$
jest typu $F_\sigma$, ponieważ
jest przeliczalną sumą singletonów
domkniętych. Nie jest typu
$G_\delta$ w $\mathbb R$.
Rzeczywiście, jako przestrzeń
z topologią względną jest
przeliczalna, a żaden jej
singleton nie ma niepustego
wnętrza. Gdyby $\mathbb Q$
była polska, twierdzenie Baire'a
z rozdziału~11 zabraniałoby
przedstawienia jej jako
przeliczalnej sumy zbiorów
nigdziegęstych. Stąd
$\mathbb Q$ nie jest polska
i z poprzedniego twierdzenia
nie jest typu $G_\delta$.

Zbiór liczb niewymiernych
$\mathbb R\setminus\mathbb Q$
jest typu $G_\delta$, a więc
jest przestrzenią polską.
Nie jest typu $F_\sigma$:
w przeciwnym razie jego
dopełnienie $\mathbb Q$
byłoby typu $G_\delta$.
\end{przyklad}

\section{Zbiory borelowskie i poziomy hierarchii}
\label{sec:topog-polskie-borel}

\begin{definicja}[Sigma-algebra borelowska]
\label{def:topog-polskie-borel}
Niech $X$ będzie przestrzenią
topologiczną. \emph{Sigma-algebrą
borelowską} $\mathcal B(X)$
nazywamy najmniejszą
$\sigma$-algebrę podzbiorów $X$
zawierającą wszystkie zbiory
otwarte. Jej elementy nazywamy
\emph{zbiorami borelowskimi}.
\end{definicja}

\begin{stwierdzenie}[Zbiory borelowskie w podprzestrzeni]
\label{prop:topog-polskie-borel-slad}
Niech $X$ będzie dowolną
przestrzenią topologiczną,
a $A\subseteq X$ podprzestrzenią.
Wtedy
\[
\mathcal B(A)=
\{A\cap B:B\in\mathcal B(X)\}.
\]
\end{stwierdzenie}
\begin{proof}
Rodzina po prawej stronie
jest $\sigma$-algebrą na $A$,
ponieważ dopełnienia względne
i przeliczalne sumy zachowują
postać $A\cap B$.
Zawiera wszystkie otwarte w $A$
zbiory $A\cap U$, gdzie $U$
jest otwarty w $X$.
Dlatego $\mathcal B(A)$ jest
w niej zawarta.

Odwrotnie, rodzina
$\mathcal C=\{B\subseteq X:
A\cap B\in\mathcal B(A)\}$
jest $\sigma$-algebrą na $X$
i zawiera wszystkie otwarte
zbiory $U\subseteq X$,
ponieważ $A\cap U$ jest
otwarty w $A$.
Stąd $\mathcal B(X)\subseteq
\mathcal C$ i każde $A\cap B$
z $B\in\mathcal B(X)$
należy do $\mathcal B(A)$.
\end{proof}

\begin{stwierdzenie}[Przeliczalne generatory sigma-algebry borelowskiej]
\label{prop:topog-polskie-borel-przeliczalna-baza}
Niech $X$ będzie przestrzenią
topologiczną spełniającą
drugi aksjomat przeliczalności
i niech $\mathcal U$ będzie
jej bazą przeliczalną.
Wówczas
\[
\mathcal B(X)=\sigma(\mathcal U),
\]
gdzie po prawej stronie oznaczamy
$\sigma$-algebrę generowaną
przez rodzinę $\mathcal U$.
W szczególności $\mathcal B(\mathbb R)$
jest generowana przez
przedziały otwarte $(a,b)$
o końcach wymiernych $a<b$.
\end{stwierdzenie}
\begin{proof}
Ponieważ wszystkie elementy
$\mathcal U$ są otwarte,
mamy $\sigma(\mathcal U)
\subseteq\mathcal B(X)$.
Każdy zbiór otwarty w $X$
jest sumą elementów $\mathcal U$;
rodzina $\mathcal U$ jest
przeliczalna, więc suma ta
jest co najwyżej przeliczalna.
Zatem każdy zbiór otwarty
należy do $\sigma(\mathcal U)$,
co daje inkluzję odwrotną.
Przedziały o końcach wymiernych
tworzą przeliczalną bazę
zwykłej topologii na $\mathbb R$.
\end{proof}

\begin{definicja}[Odwzorowanie borelowskie]
\label{def:topog-polskie-odwzorowanie-borelowskie}
Niech $X,Y$ będą przestrzeniami
topologicznymi. Odwzorowanie
$f\colon X\to Y$ nazywamy
\emph{borelowskim}, jeżeli
$f^{-1}(B)\in\mathcal B(X)$
dla każdego $B\in\mathcal B(Y)$.
\end{definicja}

\begin{stwierdzenie}[Kryterium borelowskości]
\label{prop:topog-polskie-borelowskie-kryterium}
Niech $X,Y$ będą przestrzeniami
topologicznymi, a
$f\colon X\to Y$ odwzorowaniem.
Jest ono borelowskie wtedy
i tylko wtedy, gdy
$f^{-1}(U)\in\mathcal B(X)$
dla każdego otwartego
$U\subseteq Y$.
Każde odwzorowanie ciągłe
jest borelowskie, a złożenie
dwóch odwzorowań borelowskich
jest borelowskie.
\end{stwierdzenie}
\begin{proof}
Konieczność jest oczywista.
Dla wystarczalności rozważmy
rodzinę
$\mathcal C=\{B\subseteq Y:
f^{-1}(B)\in\mathcal B(X)\}$.
Przeciwobrazy zachowują
dopełnienia i przeliczalne
sumy, więc $\mathcal C$ jest
$\sigma$-algebrą. Jeśli
zawiera wszystkie otwarte
zbiory $Y$, to zawiera
$\mathcal B(Y)$. To daje
borelowskość $f$.
Gdy $f$ jest ciągłe,
przeciwobrazy otwartych
są otwarte, więc borelowskie.
Jeśli $f\colon X\to Y$
i $g\colon Y\to Z$ są
borelowskie, to dla
$B\in\mathcal B(Z)$
mamy
$(g\circ f)^{-1}(B)
=f^{-1}(g^{-1}(B))
\in\mathcal B(X)$.
\end{proof}

\begin{definicja}[Skończone poziomy hierarchii borelowskiej]
\label{def:topog-polskie-hierarchia-skonczona}
Niech $X$ będzie przestrzenią
metryzowalną. Definiujemy
$\boldsymbol\Sigma^0_1(X)$
jako rodzinę wszystkich
zbiorów otwartych,
a $\boldsymbol\Pi^0_1(X)$
jako rodzinę wszystkich
zbiorów domkniętych.
Dla $n\ge1$ określamy
\[
\begin{aligned}
\boldsymbol\Sigma^0_{n+1}(X)
&=\left\{\bigcup_{k\ge1}A_k:
A_k\in\boldsymbol\Pi^0_n(X)\right\},\\
\boldsymbol\Pi^0_{n+1}(X)
&=\left\{\bigcap_{k\ge1}A_k:
A_k\in\boldsymbol\Sigma^0_n(X)\right\}.
\end{aligned}
\]
Są to kolejno rodziny
zbiorów otwartych i domkniętych,
zbiorów typu $F_\sigma$
i $G_\delta$, a następnie
ich dalszych przeliczalnych
sum i przecięć.
\end{definicja}

\begin{stwierdzenie}[Dualność i operacje]
\label{prop:topog-polskie-hierarchia-wlasnosci}
Niech $X$ będzie przestrzenią
metryzowalną, a $n\ge1$.
Wówczas
\[
A\in\boldsymbol\Sigma^0_n(X)
\ \Longleftrightarrow\
X\setminus A\in\boldsymbol\Pi^0_n(X).
\]
Ponadto rodzina
$\boldsymbol\Sigma^0_n(X)$
jest zamknięta na przeliczalne
sumy, a rodzina
$\boldsymbol\Pi^0_n(X)$
na przeliczalne przecięcia.
Dla każdego $n\ge1$
zachodzą inkluzje
\[
\boldsymbol\Sigma^0_n(X)
\cup\boldsymbol\Pi^0_n(X)
\subseteq
\boldsymbol\Sigma^0_{n+1}(X)
\cap\boldsymbol\Pi^0_{n+1}(X).
\]
\end{stwierdzenie}
\begin{proof}
Dualność dla $n=1$
jest definicyjna.
Jeżeli zachodzi na poziomie $n$,
to prawa de Morgana zamieniają
przeliczalną sumę zbiorów
z $\boldsymbol\Pi^0_n$
w przeliczalne przecięcie
ich dopełnień z
$\boldsymbol\Sigma^0_n$,
co daje dualność na poziomie
$n+1$.

Zamknięcie $\boldsymbol\Sigma^0_1$
na przeliczalne sumy
wynika z aksjomatów topologii.
Na poziomach $n+1$ sumę
przeliczalnie wielu przeliczalnych
sum zbiorów z
$\boldsymbol\Pi^0_n$
można zapisać jako jedną
przeliczalną sumę, indeksując
składniki przez $\mathbb N^2$.
Zamknięcie rodziny
$\boldsymbol\Pi^0_n$ na
przeliczalne przecięcia
otrzymujemy przez dopełnienia.

Lemat
\ref{lem:topog-polskie-otwarte-fsigma}
daje
$\boldsymbol\Sigma^0_1
\subseteq\boldsymbol\Sigma^0_2$
i
$\boldsymbol\Pi^0_1
\subseteq\boldsymbol\Pi^0_2$.
Ponadto każdy zbiór z
$\boldsymbol\Sigma^0_1$ jest
jednoelementowym przeliczalnym
przecięciem zbiorów otwartych,
zatem należy do
$\boldsymbol\Pi^0_2$;
analogicznie
$\boldsymbol\Pi^0_1
\subseteq\boldsymbol\Sigma^0_2$.
Przyjmijmy indukcyjnie obie
inkluzje własnych poziomów
$n\to n+1$.
Jeżeli zbiór jest przeliczalną
sumą zbiorów z
$\boldsymbol\Pi^0_n$, to
każdy składnik należy także do
$\boldsymbol\Pi^0_{n+1}$,
co daje zawieranie
$\boldsymbol\Sigma^0_{n+1}
\subseteq\boldsymbol\Sigma^0_{n+2}$.
Druga inkluzja wynika przez
dopełnienia; inkluzje krzyżowe
na każdym poziomie wynikają
z zapisania zbioru jako
przeliczalnej sumy lub
przecięcia jednego zbioru.
To kończy indukcję.
\end{proof}

\begin{uwaga}[Poziomy skończone nie wyczerpują hierarchii]
Rodziny
$\boldsymbol\Sigma^0_n$
i $\boldsymbol\Pi^0_n$,
$n\in\mathbb N$, opisują
jedynie skończone poziomy
hierarchii borelowskiej.
W ogólności pojawiają się
również zbiory borelowskie
o wyższych, przeliczalnych
rangach. Szczegółowe twierdzenia
o ścisłości hierarchii wymagają
dalszych narzędzi opisowej teorii
mnogości.
Nie należy utożsamiać całej
$\sigma$-algebry borelowskiej
z sumą jej poziomów skończonych.
\end{uwaga}

\section{Pełna hierarchia borelowska}
\label{sec:topog-polskie-hierarchia-pozaskonczona}

\begin{definicja}[Hierarchia borelowska indeksowana liczbami porządkowymi]
\label{def:topog-polskie-hierarchia-porzadkowa}
Niech $X$ będzie przestrzenią
metryzowalną. Dla każdej
niezerowej przeliczalnej
liczby porządkowej
$\alpha<\omega_1$
definiujemy przez indukcję
pozaskończoną rodziny
$\boldsymbol\Sigma^0_\alpha(X)$
i $\boldsymbol\Pi^0_\alpha(X)$.
Dla $\alpha=1$ są to
odpowiednio zbiory otwarte
i domknięte. Dla $\alpha>1$
przyjmujemy
\[
\begin{aligned}
\boldsymbol\Sigma^0_\alpha(X)
&=\left\{\bigcup_{n\ge1}A_n:
\text{ dla każdego }n\
A_n\in\boldsymbol\Pi^0_{\beta_n}(X)
\text{ dla pewnego }1\le\beta_n<\alpha
\right\},\\
\boldsymbol\Pi^0_\alpha(X)
&=\{X\setminus A:
A\in\boldsymbol\Sigma^0_\alpha(X)\}.
\end{aligned}
\]
Symbol $\omega_1$ oznacza
pierwszą nieprzeliczalną
liczbę porządkową.
Na poziomach skończonych
ta definicja zgadza się
z definicją poprzednią:
w przestrzeni metryzowalnej
rodziny $\boldsymbol\Pi^0_n$
są rosnące względem $n$,
więc suma przeliczalna
zbiorów z wcześniejszych
poziomów $\boldsymbol\Pi^0_\beta$,
$\beta<n+1$, jest taką samą
sumą zbiorów z
$\boldsymbol\Pi^0_n$.
\end{definicja}

\begin{twierdzenie}[Wyczerpanie sigma-algebry borelowskiej]
\label{thm:topog-polskie-borel-wyczerpanie}
Niech $X$ będzie przestrzenią
metryzowalną. Wówczas
\[
\mathcal B(X)=
\bigcup_{1\le\alpha<\omega_1}
\boldsymbol\Sigma^0_\alpha(X)
=
\bigcup_{1\le\alpha<\omega_1}
\boldsymbol\Pi^0_\alpha(X).
\]
\end{twierdzenie}
\begin{proof}
Każdy poziom hierarchii
jest zawarty w $\mathcal B(X)$,
co wynika przez indukcję
pozaskończoną z definicji
$\sigma$-algebry.

Oznaczmy przez $\mathcal C$
sumę poziomów
$\boldsymbol\Sigma^0_\alpha$.
Zawiera ona wszystkie zbiory
otwarte i zbiór $X$.
Jeżeli
$A\in\boldsymbol\Sigma^0_\alpha$,
to $X\setminus A$ należy do
$\boldsymbol\Pi^0_\alpha$,
a więc do
$\boldsymbol\Sigma^0_{\alpha+1}$:
można przedstawić go jako
przeliczalną sumę kopii
samego siebie, a
$\alpha<\alpha+1$.
Zatem $\mathcal C$ jest
zamknięta na dopełnienia.

Niech teraz
$A_n\in\mathcal C$,
a $\alpha_n<\omega_1$
będzie rangą taką, że
$A_n\in\boldsymbol\Sigma^0_{\alpha_n}$.
Na mocy poprzedniego argumentu
dopełnienie $A_n$ należy do
$\boldsymbol\Sigma^0_{\alpha_n+1}$,
ponieważ należy do
$\boldsymbol\Pi^0_{\alpha_n}$.
Po ponownym wzięciu dopełnienia
otrzymujemy
$A_n\in\boldsymbol\Pi^0_{\alpha_n+1}$.
Niech $\gamma$ będzie przeliczalną
liczbą porządkową większą od
każdej $\alpha_n+1$
(np. $\gamma=\sup_n(\alpha_n+2)$).
Wtedy każdy $A_n$ należy
do któregoś
$\boldsymbol\Pi^0_{\beta_n}$
z $\beta_n<\gamma$, a zatem
$\bigcup_nA_n$ należy do
$\boldsymbol\Sigma^0_\gamma$.
Ponieważ supremum przeliczalnie
wielu przeliczalnych liczb
porządkowych jest przeliczalne,
mamy $\gamma<\omega_1$.
Dowodzi to, że $\mathcal C$
jest $\sigma$-algebrą.
Jako $\sigma$-algebra
zawierająca wszystkie otwarte
musi zawierać $\mathcal B(X)$,
a więc jest jej równa.
Równość z sumą poziomów
$\boldsymbol\Pi^0_\alpha$
wynika teraz z zamknięcia
$\mathcal B(X)$ na dopełnienia
i z dualności poziomów.
\end{proof}

\begin{uwaga}[Zakres dalszego rozdziału]
W rozdziale~20 zbadamy zbiory
analityczne, obrazy ciągłe
przestrzeni polskich, własności
regularności i dalsze twierdzenia
opisowej teorii mnogości.
Hierarchia borelowska będzie
punktem odniesienia dla tych
konstrukcji; nie zastępujemy
jej samymi poziomami skończonymi.
\end{uwaga}
